\section{Introduction}

Coding agents now generate substantial code and execute multi-step workflows, but the harder problem is no longer only generation quality; it is oversight calibration. Prior work on agent oversight, appropriate reliance, and human-centered evaluation all points to the same tension: too much approval friction collapses productivity, while too much unchecked autonomy increases the risk of subtle bugs, security mistakes, and over-trust~\cite{grunde2026overseeing,bansal2021appropriate,wang2025humansmissing}. Recent field observations of professional developers reinforce that this is not a niche concern: effective use of coding agents remains highly controlled in practice, with developers deliberately steering planning, verification, and risky implementation choices rather than delegating end-to-end autonomy~\cite{huang2025dontvibe}.

Developers currently navigate this problem with a mix of markdown rule files, manual review, narrower prompts, and personal habits for when an agent should pause. These workarounds help, but they remain fragmented. Static rule files can state some explicit preferences, while repeated corrections and approval patterns reveal others only through use. Existing coding-agent workflows offer limited support for turning both kinds of signals into persistent, inspectable oversight behavior without collapsing everything into either prompt text or model-internal memory.

We present \emph{Smart Coder}, a governance-first framework for adaptive autonomy in coding agents. A formative survey with 21 experienced coding-agent users reinforced the need for milestone- and risk-based check-ins, visible reasons for changing autonomy, and less interruption on routine work. Smart Coder responds by keeping authority in a local CLI while adapting from explicit user-authored rules and observed interaction history. Its central contributions are a governed runtime in which the model remains structurally untrusted, a dual-initiator check-in design spanning both the model and the policy layer, and a trace-driven adaptation loop that changes future prompting and future oversight across sessions.
